\abstract{
We study the effect of the Large Magellanic Cloud (LMC) on the dark matter (DM) distribution in the Solar neighborhood, utilizing the Auriga magneto-hydrodynamical simulations of Milky Way (MW) analogues that have an LMC-like system. We extract the local DM velocity distribution at different times during the orbit of the LMC  around the MW in the simulations. As found in previous idealized simulations of the MW-LMC system, we find that the DM particles in the Solar neighborhood originating from the LMC analogue  dominate  the high speed tail of the local DM speed distribution. Furthermore, the native  DM particles of the MW in the Solar region are boosted to higher speeds as a result of a response to the LMC's motion.
We simulate the signals expected in near future xenon, germanium, and silicon direct detection experiments, considering  DM interactions with target nuclei or electrons. We find that the presence of the LMC causes a considerable shift in the expected direct detection exclusion limits towards smaller cross sections and DM masses, with the effect being more prominent for low mass DM. Hence, our study shows, for the first time, that the LMC's influence on the local DM distribution is significant even in fully cosmological MW analogues. 
}

